\documentclass[11pt]{article}
\usepackage[margin=2.3cm]{geometry}
\usepackage{amsmath,amssymb}
\usepackage{xr}
\externaldocument{main}
\usepackage[colorlinks=true,linkcolor=blue,urlcolor=blue]{hyperref}
\usepackage{enumitem}
\usepackage{microtype}
\setlength{\parskip}{0.45em}
\setlength{\parindent}{0pt}
\newenvironment{comment}[1]{\par\medskip\noindent\textbf{Comment~#1.}\ \itshape}{\par}
\newcommand{\resp}{\par\noindent\textbf{Response.}\ }

\title{Response to the reviewer (third round)\\[0.3em]\large A reliability-constrained, multi-objective
framework for minimum-mass design of steel frame structures under material and loading uncertainty}
\author{[Author name]}
\date{}

\begin{document}
\maketitle

We thank the reviewer for the recommendation of minor revision and for pinpointing how the evidence
should be characterised. We took both routes offered under Major comment~1. First, we quantified the
uncertainty of the residual-tail correction and now report a combined lower bound. Second, we
explicitly limited the full-model claim to the tested tail strata and the stated residual model. All
new numbers come from the seeded scripts (function \texttt{tail\_rows} in \texttt{make\_figures.py})
and are stored in \texttt{results/tail\_summary.json}. Section and table numbers refer to the revised
manuscript.

\begin{comment}{1 (three-storey full-model claim)}
The $t_3$-adjusted estimate (3.52) of the margin design has no confidence interval reflecting the
uncertainty of the fitted tail model. Either quantify it and give a conservative lower bound, or
limit the conclusion to the tested strata and the stated $t_3$ model.
\end{comment}
\resp We did both.
\begin{itemize}[itemsep=0.1em]
\item \emph{Quantified uncertainty.} The location and scale of the fitted $\ln(U_{full}/U_s)$ model
are now resampled by a nonparametric bootstrap (4000 resamples of the 550 re-analysed log-ratios
per design). Each bootstrap residual is added to a draw from the Jeffreys posterior of the sampled
strata. For the margin design (8408~kg) this gives $\hat\beta=3.52$, 95\% interval [3.41, 3.61] and
one-sided lower bound 3.43, conditional on the $t_3$ model. With the more extreme $t_2$ model the
point estimate is 3.16. Table~\ref{tab:tail} now has these columns for all ten designs, kept
separate from the sampled-strata intervals.
\item \emph{Limited claim.} Section~\ref{sec:strict} now gives two estimates of different standing:
a statistical result on the sampled strata (3.61 [3.45, 3.73]) and a model-based engineering estimate
for the complete full model (3.52, lower bound 3.43 under $t_3$). It states that the target is met with
95\% confidence on the sampled strata. For the complete model it is met only under the $t_3$
residual model, which is not a distribution-free demonstration. The abstract and Conclusion~5 use the
same distinction.
\end{itemize}

\begin{comment}{2 (drift-bias sensitivity in the headline findings)}
Qualify in the abstract and discussion that RBDO-C95 meets the targets under the baseline surrogate,
that the confidence bound addresses Monte Carlo sampling only, and that results near the
serviceability threshold are sensitive to the $B_1$--$B_2$ drift error.
\end{comment}
\resp
\begin{itemize}[itemsep=0.1em]
\item \emph{Abstract.} It now says that the RBDO-C95 designs meet both targets with 95\% confidence
``in the baseline surrogate model''. It adds that these bounds cover sampling uncertainty only and that
drift results near the target remain sensitive to the 1\% error of the $B_1$--$B_2$ drifts.
\item \emph{Methods and Results.} The same qualifier was added in Section~\ref{sec:mc}: the confidence
bound controls Monte Carlo sampling uncertainty and does not remove systematic approximation error. It
also appears in Section~\ref{sec:strict}, in the Discussion (``Statistical confidence and acceptance'')
and in Conclusion~4, which names the nine-storey design as becoming inconclusive under the 1\% drift
bias.
\item \emph{Base frame.} The RBDO-C design of the base frame is now described as meeting the drift
target only 0.007 above the threshold. A 1\% drift bias would put it below
(Section~\ref{sec:basedesign}).
\end{itemize}

\begin{comment}{3 (model dependence of the residual-tail validation)}
Explain why $t_3$ is a suitable conservative sensitivity case, and state that it is a stress-test
assumption rather than a validated tail.
\end{comment}
\resp A paragraph ``Why $t_3$, and what it is not'' was added to Section~\ref{sec:fullmc}. It states
that $t_3$ is a stress-test assumption, not a validated distribution of the extreme tail. It also
explains why it is conservative relative to the data. We compared the 5\,500 standardised
log-ratios of strata B and C of the ten designs with the two models:

\begin{center}
\begin{tabular}{lrrr}
Standardised log-ratio & Observed & Normal model & $t_3$ model \\
\hline
$>2$ & 78 & 125 & 383 \\
$>3$ & 10 & 7.4 & 159 \\
$>4$ & 0 & 0.2 & 77 \\
\end{tabular}
\end{center}

The largest observed value is 3.84. The data are close to normal up to about four standard
deviations and far lighter-tailed than $t_3$. The model is therefore pessimistic within the
observed range, while nothing is known beyond it. The paragraph ends by stating that the residual
estimates are conditional on this assumption. The $t_2$ case was added to show how much the result
depends on the tail model.

\begin{comment}{4 (scope of the novelty claim)}
Keep the full-text versus abstract distinction and restrict claims to the seven frames and the stated
assumptions.
\end{comment}
\resp The distinction between the two studies read in full and those assessed from abstracts only
($^\dagger$ in Table~\ref{tab:literature}) is retained. The contribution statement still states that we
make no claim about the latter. The scope sentence at the end of the contribution list now restricts
what is demonstrated to the seven planar frame configurations, the component-based ULS and the stated
probabilistic model.

\section*{Additional comments}
\begin{itemize}[itemsep=0.2em]
\item \emph{Abstract and conclusions.} Both now separate the baseline-surrogate result, the
sampled-strata full-model result and the residual-adjusted estimate for the three-storey frame
(Conclusion~5 is new).
\item \emph{Confidence terminology.} ``Meets the target with 95\% confidence'' is now used only where
the stated lower bound exceeds the target. The model under which this holds is always named: the
baseline surrogate model, the sampled strata, or the $t_3$ residual model.
\item \emph{Table~\ref{tab:tail}.} It is restructured into three groups: stratum C, sampled strata
only (statistical intervals), and sampled strata plus the $t$-model residual (model-conditional
intervals, with the bootstrap). The caption explains the difference.
\item \emph{Table~\ref{tab:strict}.} The caption now flags the results that become below or
inconclusive under the 1\% drift perturbation, and it states that the lower bounds address sampling
uncertainty only.
\item \emph{Table~\ref{tab:mudesign}.} The caption states that RBDO-C-MU was optimised for two frame
configurations only, and that c.o.v.\ 0.10 was not optimised.
\item \emph{Author information.} The author and affiliation placeholders remain in this version. They
must be completed by the authors before formal submission.
\item \emph{Reproducibility.} The manuscript names the commit that produced its results. A release tag
and DOI will be created on acceptance.
\end{itemize}

\end{document}
